% Exemplo: prova com a classe exam (cabeçalho, questões com pontos, itens, múltipla escolha,
% espaço para resposta e quadro de pontuação). O mesmo arquivo gera o gabarito: veja gabarito.tex.
% Compilar: latexmk -pdf prova.tex              Testado com: exam 2.7, TeX Live 2024 ou mais recente
%%BEGIN:provaclasse
% addpoints: soma os pontos das questões
\documentclass[12pt, a4paper, addpoints]{exam}
\usepackage[T1]{fontenc}
\usepackage[brazilian]{babel}
\usepackage{amsmath}
%%END:provaclasse

%%BEGIN:provaportugues
% textos da classe em português
\pointpoints{ponto}{pontos}
\renewcommand{\solutiontitle}{\noindent\textbf{Resposta:}\enspace}
% quadro de pontuação
\hqword{Questão:} \hpword{Valor:} \hsword{Nota:} \htword{Total}
\CorrectChoiceEmphasis{\bfseries}
%%END:provaportugues

%%BEGIN:provacabecalho
\pagestyle{headandfoot}
\firstpageheader{Cálculo I}{Prova 1}{\today}
\runningheader{Cálculo I}{Prova 1 (continuação)}{}
\footer{}{Página \thepage\ de \numpages}{}
\runningheadrule
%%END:provacabecalho

\begin{document}

%%BEGIN:provaidentificacao
\noindent Nome: \makebox[0.55\textwidth]{\hrulefill} \hfill
Matrícula: \makebox[3cm]{\hrulefill}

\begin{center}
  \gradetable[h][questions]
\end{center}
%%END:provaidentificacao

%%BEGIN:provaquestoes
\begin{questions}
  \question[2] Calcule a derivada de $f(x) = 3x^2 + 2x$.
  \begin{solutionorbox}[2cm] % na prova: caixa vazia; no gabarito: a resposta
    $f'(x) = 6x + 2$.
  \end{solutionorbox}

  \question Considere $g(x) = x^3$.
  \begin{parts}
    \part[1] Calcule $g'(2)$.
    \begin{solution}[1cm] $g'(x) = 3x^2$, logo $g'(2) = 12$. \end{solution}
    \part[2] Escreva a equação da reta tangente ao gráfico de $g$ em $x = 2$.
    \begin{solutionorlines}[2cm]
      $y = 12x - 16$.
    \end{solutionorlines}
  \end{parts}

  \question[1] Qual é a derivada de $\sin x$?
  \begin{choices}
    \choice $-\cos x$
    \CorrectChoice $\cos x$
    \choice $\sin x$
    \choice $-\sin x$
  \end{choices}
\end{questions}
%%END:provaquestoes

\end{document}
